\par\SourceLeaf{142}
Unter solchen Umständen tritt die Bedeutung des Amts des Vorsitzenden einer beratenden Körperschaft hervor. Er muss die einzelnen Fragen scharf voneinander trennen, jede an ihren richtigen Platz stellen und in klarer, jede Zweideutigkeit ausschließender Formulierung zur Beratung und Abstimmung bringen.

Im vorliegenden Fall oblag diese Aufgabe F. als dem damaligen Prorektor. Am 22. Juli ließ er das Reskript, das \textsc{Kritter}sche Gutachten und P.s P. M. unter sämtlichen Professoren umlaufen, mit einer Aufforderung, die ich hier vollständig in F.s eigenen Worten wiedergebe:
\begin{quote}
Sie möchten sich schriftlich dazu erklären, ob Sie P.s ganzem Vorschlag zustimmen, oder zu den einzelnen Fragen, nämlich 1) in welchem Maß die Erhöhung der Pensionen gerecht und ratsam erscheine? 2) Ob die Verlängerung der Pensionsdauer nach dem Tod der Eltern bis zum \(20^{\text{sten}}\) Lebensjahr des jüngsten Kindes gewünscht werde? 3) Ob die von der Königlichen Regierung vorgeschlagene Erhöhung der jährlichen Beiträge auf 10 Thl. C. M. genehmigt werde, ohne Bedingung oder unter der Bedingung von Nr. 2.
\end{quote}
\noindent\hspace*{\parindent}Man sieht, dass der letzte Artikel von P.s P. M. (oben S. [136]) mit keinem Wort erwähnt ist. Ich selbst habe nun zwar keinen Zweifel, dass F. ihn (vielleicht ohne die große Wichtigkeit ‘dieses Umstandes’ genügend zu würdigen) als einen bereits in seiner Nr. 1 enthaltenen Anhang betrachtet haben mag. Ebenso wenig zweifle ich aber daran, dass diese Nichterwähnung, besonders im Gegensatz zu der präzisen logischen Form, in der F. seine dritte Frage vorlegt, viel dazu beigetragen hat, jene Hauptfrage für viele, vielleicht für die meisten Abstimmenden in den Hintergrund zu rücken. Ich habe deshalb geglaubt, diese an sich geringfügigen Nebenumstände hier mitberühren zu müssen, weil dadurch die sonst so auffallende Erscheinung verständlicher wird, dass von 41 Abstimmenden auch nicht ein einziger auf eine Diskussion jener Hauptfrage eingegangen ist (wenn man nicht E.s Votum S. [139] dafür gelten lassen will); ja, dass die meisten Abstimmenden sie überhaupt nicht und eigentlich nur zwei sie überhaupt ganz unzweideutig erwähnt haben.

Diese schriftlichen Verhandlungen (vom 22. Juli bis zum 6. August) sind sehr umfangreich, und manche einzelnen Voten sehr ausführlich. Sie drehen sich aber fast ausschließlich um F.s Fragen 2 und 3, die, besonders die letzte, vielfachen Widerspruch fanden, sowie um einige neue Vorschläge, die im Lauf der Verhandlung eingebracht wurden. Dazu gehört namentlich der Vorschlag, die bisherige Freiheit aufzuheben, der Witwenkasse erst später gegen doppelte Nachzahlung der Beiträge beitreten zu können; einige Abstimmende unterstützten diesen Vorschlag, andere wiesen ihn nachdrücklich zurück. Ich werde nur einige wenige dieser Voten anführen, die mit meinem Gegenstand näher zusammenhängen.

R. erklärt sich im Allgemeinen gegen alle Veränderungen, will aber den Beschlüssen der Mehrheit zustimmen. P.s Art. 3 erwähnt er zwar überhaupt nicht, wohl aber die ihm zugrunde liegende Grundsatzfrage. Dabei tritt er der Behauptung, die P. bei der früheren Abstimmung aufgestellt hatte (oben S. [133] Z. [21]), sehr entschieden entgegen.

Er könne sich, sagt er, nicht von der Richtigkeit des Grundsatzes überzeugen, dass man bei einer Erhöhung der jetzigen Witwenpension darauf achten müsse, dass später nicht etwa die Notwendigkeit entstehe, sie zu vermindern und den künftigen Witwen weniger zu geben, als die jetzigen erhalten. Dieser Grundsatz wäre allerdings richtig, wenn die Existenz der Kasse, wie bei anderen Witwenkassen, allein auf den Beiträgen beruhte. Da aber die Professorenwitwenkasse eine Zuwendung sei, die Beiträge fast mit nichts zu veranschlagen seien und die Teilnahme an der Zuwendung für jeden Professor eine Bedingung seiner Berufung sei, gelte weit natürlicher der Grundsatz
\begin{quote}
Jede Witwe erhält die höchstmögliche Pension, die das Vermögen zu ihren Lebzeiten gestattet.
\end{quote}
\par\SourceLeaf{143}
Bei diesem Grundsatz habe niemand Grund, sich zu beklagen, ob das Vermögen steige oder falle. Sage man, die künftigen Witwen könnten weniger erhalten, wenn das Vermögen sinke, so antworte er (R.), die jetzigen Witwen könnten weniger als die künftigen erhalten, wenn das Vermögen steige; dieser letztere Fall sei wahrscheinlicher als der erstere.

R. hält also, wie man sieht, P.s Behauptung, die künftige Witwe werde geschädigt, wenn sie weniger erhalte als eine frühere erhalten hat, unausgesprochen die Erwiderung entgegen, dass man dann mit genau demselben Recht behaupten könne, die jetzige Witwe werde geschädigt, wenn sie weniger erhalte als eine künftige.

\textsc{Sch.} wiederholt den eben angeführten Grundsatz (‘Jede Witwe erhält die höchstmögliche usw.’) mit dem Zusatz: ‘H. H. R., scheint mir, hat diesen Satz zur Evidenz gebracht, worauf ich mich beziehe. \textit{Er liegt auch Nr. 3 von P.s Vorschlag zugrunde}’. Nachdem \textsc{Sch.} außerdem hervorgehoben hat, dass die Kasse eine Zuwendung, ein Teil des Gehalts und die Teilnahme daran mit einer Art von Zwang, wenn auch einem milden und gerechten, verbunden sei, fährt er fort:
\begin{quote}
Aus diesem Unterschied ergibt sich unter anderem, dass der Satz ‘wir sind den gegenwärtigen Witwen \textit{so viel} schuldig wie den künftigen und umgekehrt’ hier nicht anwendbar sei, wenn \textit{so viel} den Geldbetrag ausdrücken solle. Die Kasse hat etwas Aktienartiges, das unter der Macht der jeweiligen Umstände steht und steigt und fällt.

Erhöhung der Witwenpensionen: Ihre Höhe (\textit{wie auch die der Verminderung}) hat Herr C. R. P. durch eine unveränderliche Regel bestimmt.
\end{quote}
\noindent Ich habe diese zwei Stellen wörtlich abgeschrieben, weil daraus und namentlich aus den beiden von mir doppelt unterstrichenen ganz unwidersprechlich hervorgeht, dass \textsc{Sch.} P.s Art. Nr. 3 im Sinn der ersten Interpretation (oben S. [129]) aufgefasst hat. Mehrere andere, die danach abstimmten, wie B., H., O., R., T., haben \textsc{Sch.}s Votum zugestimmt, wenn auch ohne die einzelnen Fragen zu benennen. Andererseits erkennt man in dem [auf folgender S.] anzuführenden Votum F.s die zweite Interpretation, die auch seit 1837 in der Praxis befolgt worden ist. Ich meine daher, dass die Formulierung der Progressionsnorm, die, wie der Vergleich zeigt, unverändert in das Reskript vom 20. November 1794 übernommen worden ist, sehr wohl eine \textit{unklare} genannt werden kann.

Außer \textsc{Sch.} hat nur noch M. unsere Frage (P.s 3) ausdrücklich und unzweideutig erwähnt. Er sagt aber nichts weiter darüber, als dass das von P. empfohlene \textit{Festsetzen} einer Norm für das künftige Steigen der Pension besonders wichtig scheine. Man könnte sagen, genau genommen liege darin nur eine Billigung des Zwecks, aber noch nicht bestimmt eine Billigung des von P. vorgeschlagenen Mittels und ebenso wenig eine Erklärung, in welchem Sinn M. Letzteres aufgefasst habe. Ohne darauf jedoch Gewicht zu legen, will ich nicht unerwähnt lassen, dass \textsc{Meiners} in seinem bekannten Werk über die Verfassung und Verwaltung deutscher Universitäten S. 95 die bestehende Progressionsnorm ebenso wie \textsc{Brandes} (siehe oben S. [129]) mit den Worten anführt: \textit{so lange} die Zahl der Witwen nicht über 15 hinausgehe.

Von P. selbst liegt zu unserer Frage nichts weiter vor als der oben S. [136] mitgeteilte Art. 3. Darf ich noch einen Augenblick bei der Frage verweilen, in welchem Sinn P. selbst ihn verstanden hat, so mache ich darauf aufmerksam, dass die Worte ‘Es ist klar, dass man dies tun kann’ nur auf zwei Arten ausgelegt werden können. Entweder war P.s Meinung, eine solche Erhöhung
\begin{quote}
solle nur so lange gelten, wie die Zahl 15 noch nicht überschritten sei, danach aber entweder wieder wegfallen oder in angemessener Weise abgeändert werden
\end{quote}
\noindent oder P. hat das Überschreiten der Zahl 15 für unmöglich, wenigstens aber für so sehr unwahrscheinlich gehalten, dass es ganz unnötig sei, einen solchen Fall zu berücksichtigen.
